% Includable manuscript fragment. Requires amsmath, amssymb, amsthm.
% Requires theorem and lemma environments; no project-specific macros.
% Canonical mathematical calculus only. Raw JSON alias soundness is separate.
\section{A deterministic introduction rule selected by the witness}
\label{sec:deterministic-witness}

The independent case/progress choice can be removed from paired
introduction without changing the native consequence relation.  The
replacement selects progress precisely when the matrix is a literal
numerical equation and the supplied witness is the literal zero term.
Every other introduction uses the zero case.  This is a restriction on
finite syntax, not a test of the value denoted by a term.

\subsection{The two precisely specified calculi}

Numerical terms are finite expressions
\[
 t::=x\mid0\mid St\mid I(n,a;v.h),
\]
where only $v$ in $h$ is bound.  Every update is fixed finite syntax with
passive parameters, and all displayed substitutions avoid capture.
For the executable interface we use the conservative closure extension of
the Eq/$D$ language:
\[
 A::=(s=t)\mid D_{x,y}A\mid Q_yA.
\]
The primary Eq/$D$ fragment is recovered by replacing $Q_yB$ with
$D_{u,y}B$ for an unused $u$. The two expressions are interderivable:
paired elimination gives the closure in one direction, and case
introduction with the unused witness $S0$, using its base and QE,
gives the other. The inherited
logical bridge explains this conservative notation.
The ordinary reading of a closure is
$B(0)\to\forall y\,B(y)$; the ordinary reading of $D_{x,y}B$ is
$\exists x\,(B(x,0)\to\forall y\,B(x,y))$.  Neither reading is
changed by the construction below.  Contexts are finite reusable lists;
auxiliary closures may occur as local resources.

The two calculi have the same computation axioms.  The proof works with
the unrestricted iterator equations, and also with the smaller
certified-zero inventory
\begin{align*}
 I(0,a;v.h)&=a
 &&\text{when }a\not\equiv_{\rm syn}0,\\
 I(Sn,a;v.h)&=h[I(n,a;v.h)/v]
 &&\text{when }n,a\not\equiv_{\rm syn}0.
\end{align*}
Here $\equiv_{\rm syn}$ denotes literal syntactic identity.  Computation
axioms may be leaves in every context.  Equivalently, a checker may supply
such an axiom as a local assumption and subsequently use hypothesis
closure; the cut lemma below justifies the equivalence.

The common logical rules are exactly hypothesis use, equality transport,
closure specialization, and paired elimination:
\[
 \frac{}{\Gamma\vdash A}\ (A\in\Gamma),
 \qquad
 \frac{\Gamma\vdash s=t\qquad\Gamma\vdash F(s)}
      {\Gamma\vdash F(t)},
\]
\[
 \frac{\Gamma\vdash Q_yB(y)\qquad\Gamma\vdash B(0)}
      {\Gamma\vdash B(r)},
 \qquad
 \frac{\Gamma\vdash D_{x,y}B(x,y)\qquad
       \Gamma,Q_yB(b,y)\vdash A}
      {\Gamma\vdash A}.
\]
The numerical parameter $b$ in paired elimination is fresh for the
retained context, conclusion, and passive matrix parameters.  The
transport motive $F$ is any fixed native formula template, including
one whose hole lies beneath a binder; binders are first renamed away
from the substituted terms.

Let $\mathcal C$ additionally have the two inherited introductions
\[
 \frac{\Gamma,B(t,0)\vdash B(t,Sz)}
      {\Gamma\vdash D_{x,y}B(x,y)}
 \quad\text{(case)},
 \qquad
 \frac{\Gamma,E(t,z)\vdash E(t,Sz)}
      {\Gamma\vdash D_{x,y}E(x,y)}
 \quad\text{(progress)},
 \tag{DI$_{\rm old}$}
\]
where $E$ must be a literal equation.  In both rules $z$ is fresh for
$\Gamma$, the witness $t$, and every passive matrix parameter.

For the new calculus $\mathcal C_{\rm w}$ define, on fixed syntax,
\[
 \delta_{B,t}(z)=
 \begin{cases}
 z,& B\text{ is a literal equation and }t\equiv_{\rm syn}0,\\
 0,&\text{otherwise}.
 \end{cases}
\]
Its only paired-introduction pattern is
\[
 \frac{\Gamma,B(t,\delta_{B,t}(z))\vdash B(t,Sz)}
      {\Gamma\vdash D_{x,y}B(x,y)}.
 \tag{DI$_{\rm w}$}\label{eq:det-witness-intro}
\]
It has the same freshness conditions.  In particular, a nonliteral
term denoting zero is still a case witness.  Compound matrices always
use cases.  No additional induction, substitution, reflexivity, symmetry,
or matrix-congruence rule is assumed.

\subsection{Structural facts established before the simulation}

\begin{lemma}[Renaming, weakening, and ordinary-assumption cut]
\label{lem:det-witness-cut}
In $\mathcal C_{\rm w}$, injective renaming of numerical variables,
weakening of reusable contexts, and ordinary-assumption cut are
admissible.  The latter has the form
\[
 \frac{\Gamma\vdash A\qquad\Gamma,A\vdash H}
      {\Gamma\vdash H},
\]
where $A$ is an object formula.  The same conclusions hold with additional
retained local closure clauses.
\end{lemma}
\begin{proof}
Variable renaming leaves both tests in
\eqref{eq:det-witness-intro} unchanged: a literal equation remains a
literal equation, and renaming never turns a nonzero syntactic term
into the literal constant $0$, or conversely.  It also preserves the
syntactic conditions on computation axioms.  Structural induction on
a derivation therefore proves renaming admissible.  Injections on the
finitely many variables occurring in a proof suffice; fresh variables
can always be chosen outside another fixed finite derivation.

For weakening, first rename local eigenparameters away from the added
context, then retain that context throughout the original proof.
Every inference is still an instance of the same rule, with the same
witness syntax.  Exchange and contraction are harmless because contexts
are reusable lists and hypothesis closure is membership.

For cut, induct on the derivation of $\Gamma,A\vdash H$.  At a use
of the cut assumption, insert the given derivation of $A$, weakened by
any local assumptions.  Other hypotheses and computation leaves are
unchanged.  Transport and closure specialization commute with this
replacement.  At a paired elimination or introduction, first rename
its eigenparameter away from the finite proof being inserted.  Apply
the induction hypothesis to its premise derivation or derivations and
repeat the original inference.  This operation substitutes derivations
for hypothesis uses; it never substitutes numerical terms into the
witness field.  Consequently the literal-zero test and the
literal-equation test are preserved at every introduction.  This proves
cut without assuming admissibility of arbitrary numerical substitution.
\end{proof}

\begin{lemma}[A closed nonliteral representative of zero]
\label{lem:det-witness-zero}
Both calculi prove $C=0$ and $0=C$, where
\[
 C:=I(S(S0),S0;v.0).
\]
These proofs use only permitted computation axioms and equality
transport, even with the restricted computation inventory above.
\end{lemma}
\begin{proof}
The successor computation with predecessor $S0$ and initializer $S0$
gives $C=0$ directly: both supplied terms are nonliteral zero, and the
constant update returns $0$.

To obtain the reverse equation without presupposing symmetry or
reflexivity, put $K:=I(0,C;v.0)$.  The permitted base computation gives
$K=C$, since its initializer $C$ is not the literal zero term.  Apply
equality transport along $K=C$ to a second copy of $K=C$, with motive
$F(w):=(w=C)$.  The conclusion is $C=C$.  Now transport this conclusion
along $C=0$, with the same motive, obtaining $0=C$.

Thus neither a base computation with literal-zero initializer nor an
introduction rule is used in the construction.  If desired, $0=0$
follows separately by transporting $C=0$ along itself with motive
$w=0$; it is not needed as a hidden premise of the displayed proof.
\end{proof}

\begin{lemma}[Computation repair independent of introduction]
\label{lem:det-witness-computation-repair}
Every unrestricted iterator computation equation is derivable from the
restricted computation inventory by equality transport alone.  No
paired introduction, ordinary-assumption cut, or admissible numerical
substitution is needed for this assertion.
\end{lemma}
\begin{proof}
The theorem $C=0$ was obtained directly in
Lemma~\ref{lem:det-witness-zero}.  A base equation with initializer
literal zero is repaired using the fixed motive
\[
 F(w):=\bigl(I(0,w;v.h)=w\bigr).
\]
Its instance $F(C)$ is a permitted computation axiom.  Transport along
$C=0$ gives $F(0)$, the missing computation.

For a successor computation, replace each forbidden literal-zero
supplied field by the same fresh placeholder $w$.  When only the
initializer is forbidden, use
\[
 F(w):=\bigl(I(Sn,w;v.h)=h[I(n,w;v.h)/v]\bigr).
\]
When only the predecessor is forbidden, use
\[
 F(w):=\bigl(I(Sw,a;v.h)=h[I(w,a;v.h)/v]\bigr).
\]
When both are forbidden, use
\[
 F(w):=\bigl(I(Sw,w;v.h)=h[I(w,w;v.h)/v]\bigr).
\]
In each case $F(C)$ is a permitted successor computation.  Transport
along $C=0$ gives the exact desired equation $F(0)$.  Choose $w$ fresh
for $h$ and all passive parameters, and rename the bound state variable
before forming each motive.  Hence only the displayed argument positions
are changed; unrelated zero occurrences and passive parameters stay fixed.
This is a construction of individual proofs by the primitive transport
rule, not an assumed operation substituting terms into whole proofs.
\end{proof}

\subsection{Exact preservation of native sequents}

\begin{theorem}[Witness-selected introduction]
\label{thm:det-witness-equivalence}
For every finite native context $\Gamma$ and object formula $A$,
\[
 \Gamma\vdash_{\mathcal C} A
 \quad\Longleftrightarrow\quad
 \Gamma\vdash_{\mathcal C_{\rm w}} A.
\]
Both translations are effective transformations of finite derivations.
The terms, formulas, contexts, and final sequent are unchanged.  In
particular the new rule has the original closure interpretation, not a
new progressive interpretation of atomic paired formulas.
\end{theorem}
\begin{proof}
Every application of DI$_{\rm w}$ is literally an old progress
application or an old case application.  The common rules are
unchanged.  Thus a new derivation is already an old derivation, proving
one inclusion without any translation of formulas.

For the reverse inclusion, induct on the old derivation, recursively
transforming its premise derivations first.  Common rules are repeated
unchanged.  Four introduction situations exhaust the remaining cases.

\paragraph{1. Old case, with compound matrix or nonliteral-zero witness.}
If $B$ is compound, or if $t\not\equiv_{\rm syn}0$, then
$\delta_{B,t}(z)=0$.  The transformed premise is exactly a premise of
DI$_{\rm w}$, so repeat the introduction.

\paragraph{2. Old case, with literal equation matrix and literal-zero
witness.}
Write the transformed premise as
\[
 \pi:\qquad\Gamma,E(0,0)\vdash E(0,Sz).
\]
Choose its eigenparameter $z$ fresh for all retained data.  Under the
new context $\Gamma,E(C,0)$, transport the hypothesis $E(C,0)$ along
the closed theorem $C=0$, using the fixed motive $E(w,0)$.  This gives
$E(0,0)$.  Lemma~\ref{lem:det-witness-cut} inserts this derivation
at the uses of the assumption $E(0,0)$ in $\pi$, giving
\[
                 \Gamma,E(C,0)\vdash E(0,Sz).
\]
Transport the resulting conclusion along $0=C$, now with motive
$E(w,Sz)$.  The conclusion is
\[
                 \Gamma,E(C,0)\vdash E(C,Sz).
\]
The term $C$ is closed and syntactically nonzero.  Consequently
DI$_{\rm w}$ with witness $C$ selects the case premise just obtained
and concludes the original formula $D_{x,y}E(x,y)$.  The conclusion's
syntax does not record which witness was used.

This argument changes the retained ordinary assumption by proved
transport and cut.  It does not apply a numerical substitution to
$\pi$, so it cannot accidentally change the selected mode of an
introduction inside that derivation.

\paragraph{3. Old progress, with literal-zero witness.}
Its matrix is a literal equation and its witness is the literal zero
term.  Therefore DI$_{\rm w}$ selects progress and repeats the old
inference using the recursively transformed premise.

\paragraph{4. Old progress, with a nonliteral-zero witness.}
Write its transformed premise as
\[
 \pi:\qquad\Gamma,E(t,z)\vdash E(t,Sz),
 \qquad t\not\equiv_{\rm syn}0.
\]
Choose a new variable $u$ absent from $\Gamma,t$ and the entire matrix.
Consider the equation template $\widetilde E(u,y):=E(t,y)$, whose
$u$ coordinate is unused.  DI$_{\rm w}$ with witness $0$ selects
progress and uses $\pi$ directly to give
\[
                    \Gamma\vdash D_{u,y}E(t,y).
\]
Paired elimination of this formula introduces the local closure
$Q_yE(t,y)$.  Its fresh existential parameter is absent from this
closure, because $u$ was unused.

Under $\Gamma,Q_yE(t,y),E(t,0)$, hypothesis closure and QE give
$E(t,Sz')$, for a newly chosen parameter $z'$ fresh for all retained
context and passive data.  Since $t$ is syntactically nonzero,
DI$_{\rm w}$ with witness $t$ selects cases and proves
\[
          \Gamma,Q_yE(t,y)\vdash D_{x,y}E(x,y).
\]
The paired elimination now discharges its local closure and gives
exactly the original conclusion.  The original premise proof $\pi$
has been inserted once.  Neither its witness nor its numerical
parameters have been substituted.

All fresh parameters used in these four constructions can be chosen
outside the finitely many variables of the transformed premises.
Renaming preserves the deterministic guard by
Lemma~\ref{lem:det-witness-cut}.  This completes the induction on the
old derivation and proves the second inclusion.
\end{proof}

\paragraph{Substitution and arithmetic repairs are consequences, not
premises.}
Numerical substitution can change a witness from a variable into the
literal constant $0$, and can also create a forbidden literal-zero
initializer in a restricted computation.  It is therefore not valid to
claim that arbitrary numerical substitution preserves every primitive
new rule instance.  Once Theorem~\ref{thm:det-witness-equivalence} is
proved, the inherited admissible substitution transformation may be
performed in the old calculus and its result translated back.  Likewise,
any separately established repair of the restricted computation axioms
transfers through the identity on native sequents.  The proof above used
neither assertion circularly.

\paragraph{Consequences and limits.}
The theorem transfers the inherited whole-domain bi-interpretation
with full first-order intuitionistic PRA, and its equiconsistency, to
$\mathcal C_{\rm w}$: compose the inherited proof translations with
the two identity-sequent transformations just proved.  All formula
round trips and numerical interpretation composites remain those of
the old calculus.  The new contribution is a restriction of primitive
proof choices, rather than a different numerical algebra or a new
meaning of intuitionistic logic.  No absolute minimum, lower-bound
optimality, or universal bound on the shortest certificates follows.
A concrete checker-size claim additionally requires the complete
implementation, its raw-input soundness argument, and measured source.

\paragraph{Attribution and the scope of the novelty search.}
The unused-witness progress construction is already present in the
inherited Choice-DI proofs, and the closed term $C$ and its direct zero
certificate are already present in the certified-zero draft.  Existing
progressive-binder proposals also use deterministic introduction, but
change the atomic closure interpretation and the corresponding
elimination rules.  The construction here combines the first two
inherited mechanisms to select the old case/progress rules by witness
syntax while preserving every old native sequent.  Searches of the
seventeen current extracted archives and their included text/code
references did not locate this exact restriction and simulation.
That bounded search is not a claim of literature-wide priority.

\paragraph{A general form of the normalization.}
The argument does not depend on the numerical value zero having a special
induction-theoretic role in the witness coordinate.  If $p$ is any fixed
closed term and there is a syntactically different closed term $q$ with
proofs of $p=q$ and $q=p$ in the common fragment, progress may instead be
selected exactly at witness $p$.  Repair a forbidden case witness $p$
using $q$, and introduce every auxiliary unused-witness progress formula
with witness $p$.  The rest of the simulation is identical.  Zero is
chosen here because its syntax is the cheapest recognizable witness and
its nonliteral representative $C$ has a direct computation certificate.
The induction coordinate itself still begins at zero; only the independent
existential witness coordinate is being normalized.
